\section{Гомоморфизмы простых расширений}

Пусть $F$~--- поле, а $E$ и $E'$~--- поля, содержащие $F$. \emph{$F$-гомоморфизм} $\varphi\colon E\to E'$~--- это
гомоморфизм полей, тождественный на $F$. Он переводит <<многочлен от элементов $E$>>
\[
\sum a_{i_1\cdots i_m}\alpha_1^{i_1}\cdots\alpha_m^{i_m}\qquad(a_{i_1\cdots i_m}\in F,\ \alpha_i\in E)
\quad\text{в}\quad
\sum a_{i_1\cdots i_m}\varphi(\alpha_1)^{i_1}\cdots\varphi(\alpha_m)^{i_m},
\]
т.\,е.\ \emph{$F$-гомоморфизм полностью определяется образами образующих}. Каждый гомоморфизм полей инъективен
(раздел <<Поля>> главы~1), поэтому $F$-гомоморфизм $E\to E'$~--- это, в частности, инъективное $F$-линейное отображение
векторных пространств над $F$. Следовательно, \emph{если $[E:F]=[E':F]<\infty$, то каждый $F$-гомоморфизм $E\to E'$~---
изоморфизм}: инъективное линейное отображение пространств одной конечной размерности сюръективно.

\begin{proposition}[2.1]
Пусть $F(\alpha)$~--- простое расширение поля $F$, а $\Om$~--- ещё одно расширение $F$.
\begin{enumerate}[label=\textup{(\alph*)}]
\item Пусть $\alpha$ трансцендентен над $F$. Для каждого $F$-гомоморфизма $\varphi\colon F(\alpha)\to\Om$ элемент
$\varphi(\alpha)$ трансцендентен над $F$, и сопоставление $\varphi\mapsto\varphi(\alpha)$ задаёт взаимно однозначное
соответствие
\[
\{F\text{-гомоморфизмы }F(\alpha)\to\Om\}\ \longleftrightarrow\ \{\text{элементы }\Om,\text{ трансцендентные над }F\}.
\]
\item Пусть $\alpha$ алгебраичен над $F$, и $f$~--- его минимальный многочлен. Для каждого $F$-гомоморфизма
$\varphi\colon F[\alpha]\to\Om$ элемент $\varphi(\alpha)$~--- корень $f$ в $\Om$, и сопоставление $\varphi\mapsto\varphi(\alpha)$
задаёт взаимно однозначное соответствие
\[
\{F\text{-гомоморфизмы }F[\alpha]\to\Om\}\ \longleftrightarrow\ \{\text{корни }f\text{ в }\Om\}.
\]
В частности, число таких гомоморфизмов равно числу \emph{различных} корней $f$ в $\Om$.
\end{enumerate}
\end{proposition}

\begin{proof}
В обоих случаях отображение $\varphi\mapsto\varphi(\alpha)$ инъективно: $F$-гомоморфизм определяется образом
образующей $\alpha$ (на $F[\alpha]$~--- по формуле $\sum a_i\alpha^i\mapsto\sum a_i\varphi(\alpha)^i$, а на дробях
$g(\alpha)/h(\alpha)$~--- как частное образов). Остаётся описать множество возможных образов.

(a) Трансцендентность $\alpha$ означает, что $X\mapsto\alpha$ задаёт изоморфизм $F[X]\to F[\alpha]$ (глава~1, случай~1
перед \rf{1.27c}). Поэтому по универсальному свойству кольца многочленов (\rf{1.5}) для каждого $\gamma\in\Om$ существует
единственный $F$-гомоморфизм $\varphi\colon F[\alpha]\to\Om$ с $\varphi(\alpha)=\gamma$, а именно $g(\alpha)\mapsto g(\gamma)$.
Он продолжается до поля частных $F(\alpha)$ кольца $F[\alpha]$ (и тогда единственным образом: $g(\alpha)/h(\alpha)\mapsto
g(\gamma)/h(\gamma)$) в том и только в том случае, когда ненулевые элементы $F[\alpha]$ переходят в ненулевые, т.\,е.\
когда $h(\gamma)\ne0$ для каждого ненулевого $h\in F[X]$. Последнее и означает, что $\gamma$ трансцендентен над $F$.
Итак, имеются взаимно однозначные соответствия между $F$-гомоморфизмами $F(\alpha)\to\Om$, $F$-гомоморфизмами
$F[\alpha]\to\Om$ с трансцендентным образом $\alpha$ и трансцендентными над $F$ элементами $\Om$.

(b) Пусть $f=\sum a_iX^i$ и $\varphi\colon F[\alpha]\to\Om$~--- $F$-гомоморфизм. Применяя $\varphi$ к равенству
$\sum a_i\alpha^i=0$ и пользуясь тем, что $\varphi(a_i)=a_i$, получаем $\sum a_i\varphi(\alpha)^i=0$, т.\,е.\
$\varphi(\alpha)$~--- корень $f$ в $\Om$. Обратно, пусть $\gamma\in\Om$~--- корень $f$. Гомоморфизм подстановки
$F[X]\to\Om$, $g\mapsto g(\gamma)$, переводит $f$ в нуль, поэтому его ядро содержит $(f)$, и он пропускается через
$F[X]/(f)$. Композиция с обратным к каноническому изоморфизму $F[X]/(f)\to F[\alpha]$, $X+(f)\mapsto\alpha$ (\rf{1.27c}),
даёт $F$-гомоморфизм $F[\alpha]\to\Om$, переводящий $\alpha$ в $\gamma$.
\end{proof}

\par\noindent\begin{minipage}{\linewidth}
\begin{center}
\begin{tikzcd}[column sep=large,row sep=large]
F[X]\arrow[r,cold,"g\,\mapsto\,g(\gamma)"]\arrow[d,two heads] & \Om\\
F[X]/(f)\arrow[ur,cnew,"x\,\mapsto\,\gamma"'] & F[\alpha]\arrow[l,"\sim"',"\alpha\,\mapsto\,x"]\arrow[u,cnew,dashed,"\alpha\,\mapsto\,\gamma"']
\end{tikzcd}
\end{center}
\figcap{Построение $F$-гомоморфизма $F[\alpha]\to\Om$ по корню $\gamma$ минимального многочлена $f$: подстановка $\gamma$
обнуляет $f$ и потому пропускается через $F[X]/(f)\cong F[\alpha]$ (здесь $x=X+(f)$~--- класс $X$).}
\end{minipage}\par\medskip

\begin{remark}[2.1a]
(1) В (b) корни считаются \emph{без учёта кратностей}: гомоморфизмов столько, сколько различных элементов
$\gamma\in\Om$ с $f(\gamma)=0$. Так как у $f$ не больше $\deg f$ корней (\rf{1.7}), число $F$-гомоморфизмов $F[\alpha]\to\Om$
не превосходит $\deg f=[F[\alpha]:F]$; равенство означает, что $f$ имеет в $\Om$ ровно $\deg f$ различных корней, т.\,е.\
раскладывается в $\Om[X]$ на различные линейные множители. Гомоморфизмов может и не быть вовсе: $\Q[i]\to\R$.

(2) В (a) для кольца $F[\alpha]$ годится \emph{любой} образ $\gamma\in\Om$, а для поля $F(\alpha)$~--- только
трансцендентный. Например, если $\Om$ алгебраично над $F$ (скажем, $\Om=F$), то $F$-гомоморфизмов $F(\alpha)\to\Om$ нет
совсем, хотя гомоморфизмов $F[\alpha]\to\Om$ много (<<подстановка $\alpha=\gamma$>>; она не продолжается на дробь
$1/h(\alpha)$, где $h$~--- минимальный многочлен $\gamma$ над $F$; при $\Om=F$ это $1/(\alpha-\gamma)$).
\end{remark}

\begin{example}[2.2]
Пусть $\Q[\alpha]$~--- простое алгебраическое расширение $\Q$ и $f$~--- минимальный многочлен $\alpha$. Позже мы увидим,
что $f$ имеет в $\CC$ ровно $\deg f$ различных корней (неприводимый многочлен над полем характеристики нуль не имеет
кратных корней, \rf{2.20}; а $\CC$ алгебраически замкнуто, \rf{5.6}). Поэтому существует ровно $[\Q[\alpha]:\Q]$
различных $\Q$-гомоморфизмов $\Q[\alpha]\to\CC$, и каждый переводит $\alpha$ в корень $f$ в $\CC$. Например, для
$\alpha=\sqrt[3]2$ это три гомоморфизма $\sqrt[3]2\mapsto\sqrt[3]2,\ \omega\sqrt[3]2,\ \omega^2\sqrt[3]2$
($\omega=e^{2\pi i/3}$), а в $\R$ поле $\Q[\sqrt[3]2]$ вкладывается лишь одним способом.
\end{example}

\begin{example}[2.3]
Пусть $F$~--- поле характеристики $p\ne0$ и $a\in F\smallsetminus F^p$ (т.\,е.\ $a$ не является $p$-й степенью в $F$).
\emph{Тогда многочлен $X^p-a$ неприводим.} Пусть $F[\alpha]$~--- соответствующее поле корня. Если $\beta$~--- корень
$X^p-a$ в каком-либо расширении $\Om$ поля $F$, то $X^p-a=(X-\beta)^p$ в $\Om[X]$, и поэтому существует ровно один
$F$-гомоморфизм $F[\alpha]\to\Om$ (переводящий $\alpha$ в $\beta$).

\emph{Доказательство.} Сначала найдём расширение, в котором у $X^p-a$ есть корень; алгебраическое замыкание для этого
не нужно. Возьмём унитарный неприводимый делитель $g$ многочлена $X^p-a$ (он есть, так как $F[X]$ факториально) и
положим $E=F[X]/(g)$; по \rf{1.25} это поле, и класс $\beta$ переменной $X$~--- корень $g$, а значит, и $X^p-a$. В
кольце $E[X]$ характеристики $p$ возведение в степень $p$~--- гомоморфизм (эндоморфизм Фробениуса, глава~1), поэтому
\[
X^p-a=X^p-\beta^p=(X-\beta)^p .
\]
Пусть теперь $h\in F[X]$~--- \emph{любой} унитарный неприводимый делитель $X^p-a$. Он делит $(X-\beta)^p$ и в $E[X]$, а
по однозначности разложения в $E[X]$ унитарные делители $(X-\beta)^p$~--- это $(X-\beta)^i$; итак, $h=(X-\beta)^i$ для
некоторого $1\le i\le p$. Коэффициент при $X^{i-1}$ в $(X-\beta)^i$ равен $-i\beta$ и лежит в $F$. Если $i<p$, то
$i\cdot1_F\ne0$ обратимо в $F$, откуда $\beta\in F$ и $a=\beta^p\in F^p$~--- противоречие. Значит, $i=p$ и $h=X^p-a$, т.\,е.\
$X^p-a$ неприводим.

Пусть теперь $\beta\in\Om$~--- корень $X^p-a$ в каком-либо расширении $\Om$. Те же рассуждения в $\Om[X]$ дают
$X^p-a=(X-\beta)^p$, так что $\beta$~--- единственный корень $X^p-a$ в $\Om$, и по \rf{2.1}(b) существует ровно один
$F$-гомоморфизм $F[\alpha]\to\Om$.
\end{example}

Здесь $[F[\alpha]:F]=p$, а $F$-гомоморфизм в любое расширение, содержащее корень $X^p-a$,~--- \emph{один} (если корня
нет, например при $\Om=F$, гомоморфизмов нет совсем): оценка \rf{2.1a}(1) может быть строгой даже
тогда, когда многочлен раскладывается в $\Om$ на линейные множители. Такие элементы $a$ существуют не всегда. Если $F$
конечно, то эндоморфизм Фробениуса $x\mapsto x^p$ инъективен и, значит, сюръективен, т.\,е.\ $F^p=F$; то же верно для
алгебраически замкнутого $F$. Простейший пример поля с $F^p\ne F$~--- поле рациональных функций $F=\F_p(T)$ и $a=T$
(\rf{1.4a}). К этому примеру мы вернёмся в разделе о совершенных полях.

\begin{remark}[2.3a]
(1) \emph{Неприводимость в \rf{2.3} можно получить и из свободного члена.} Пусть, как в доказательстве, $h=(X-\beta)^i$,
$1\le i\le p$,~--- неприводимый делитель $X^p-a$ в $F[X]$. Его свободный член $(-\beta)^i$ лежит в $F$, т.\,е.\
$\beta^i\in F$; кроме того, $\beta^p=a\in F$ и $\beta\ne0$ (ведь $a\ne0$, так как $0=0^p\in F^p$). Если $i<p$, то
$\gcd(i,p)=1$, и найдутся $u,v\in\Z$ с $ui+vp=1$; тогда $\beta=(\beta^i)^u(\beta^p)^v\in F$ и $a=\beta^p\in F^p$~---
противоречие. Значит, $i=p$ и $h=X^p-a$.

(2) \emph{Пусть $n\ge1$. Многочлен $X^{p^n}-a$ неприводим в $F[X]$ тогда и только тогда, когда $a\notin F^p$.} Если
$a=c^p$, $c\in F$, то $X^{p^n}-a=\bigl(X^{p^{n-1}}-c\bigr)^p$ (Фробениус в $F[X]$). Пусть $a\notin F^p$, $h$~--- унитарный
неприводимый делитель $X^{p^n}-a$ и $\beta$~--- класс $X$ в поле $E=F[X]/(h)$. Как в \rf{2.3},
$X^{p^n}-a=X^{p^n}-\beta^{p^n}=(X-\beta)^{p^n}$ в $E[X]$, поэтому $h=(X-\beta)^i$, $1\le i\le p^n$. Запишем $i=p^sm$ с
$p\nmid m$; тогда $s\le n$ и
\[
h=\bigl((X-\beta)^{p^s}\bigr)^m=\bigl(X^{p^s}-b\bigr)^m=X^{p^sm}-mb\,X^{p^s(m-1)}+\dots,\qquad b=\beta^{p^s}.
\]
Коэффициент $-mb$ лежит в $F$, а $m\cdot1_F\ne0$; значит, $b\in F$. Если $s<n$, то
$a=\beta^{p^n}=b^{p^{n-s}}=\bigl(b^{p^{n-s-1}}\bigr)^p\in F^p$~--- противоречие. Итак, $s=n$, откуда $m=1$ (так как $i\le p^n$) и
$h=(X-\beta)^{p^n}=X^{p^n}-a$. Например, $X^{p^n}-T$ неприводим над $\F_p(T)$, и в поле его корня у него единственный корень
кратности $p^n$.
\end{remark}

Нам понадобится небольшое обобщение предложения \rf{2.1}: вместо $F$-гомоморфизмов рассматриваются продолжения
заданного гомоморфизма $\varphi_0\colon F\to\Om$. Для $f=\sum a_iX^i\in F[X]$ обозначим через $\varphi_0f=\sum\varphi_0(a_i)X^i\in\Om[X]$
многочлен, полученный применением $\varphi_0$ к коэффициентам. \emph{Продолжение} $\varphi_0$ на $F(\alpha)$~--- это
гомоморфизм $\varphi\colon F(\alpha)\to\Om$, ограничение которого на $F$ равно $\varphi_0$. Отображение $f\mapsto\varphi_0f$~---
гомоморфизм колец $F[X]\to\Om[X]$, а $\varphi_0\colon F\to\varphi_0(F)$~--- изоморфизм полей.

\begin{proposition}[2.4]
Пусть $F(\alpha)$~--- простое расширение $F$, а $\varphi_0\colon F\to\Om$~--- гомоморфизм в поле $\Om$.
\begin{enumerate}[label=\textup{(\alph*)}]
\item Если $\alpha$ трансцендентен над $F$, то $\varphi\mapsto\varphi(\alpha)$ задаёт взаимно однозначное соответствие
\[
\{\varphi\colon F(\alpha)\to\Om\mid\varphi|_F=\varphi_0\}\ \longleftrightarrow\
\{\gamma\in\Om\text{ трансцендентные над }\varphi_0(F)\}.
\]
\item Если $\alpha$ алгебраичен над $F$ с минимальным многочленом $f$, то $\varphi\mapsto\varphi(\alpha)$ задаёт взаимно
однозначное соответствие
\[
\{\varphi\colon F[\alpha]\to\Om\mid\varphi|_F=\varphi_0\}\ \longleftrightarrow\
\{\text{корни }\varphi_0f\text{ в }\Om\}.
\]
В частности, число таких продолжений равно числу различных корней $\varphi_0f$ в $\Om$ и не превосходит
$\deg f=[F[\alpha]:F]$.
\end{enumerate}
\end{proposition}

\begin{proof}
Доказательство повторяет доказательство \rf{2.1}; приведём его полностью. Инъективность $\varphi\mapsto\varphi(\alpha)$
та же: продолжение $\varphi$ обязано переводить $\sum a_i\alpha^i$ в $\sum\varphi_0(a_i)\varphi(\alpha)^i$, а дроби~--- в частные
образов.

Для $\gamma\in\Om$ рассмотрим гомоморфизм колец
\[
\theta_\gamma\colon F[X]\to\Om,\qquad g\mapsto(\varphi_0g)(\gamma)
\]
(композиция $g\mapsto\varphi_0g$ и подстановки $\gamma$). Он совпадает с $\varphi_0$ на $F$ и переводит $X$ в $\gamma$.

(a) Так как $F[X]\to F[\alpha]$, $X\mapsto\alpha$,~--- изоморфизм, гомоморфизмы $F[\alpha]\to\Om$, продолжающие $\varphi_0$,~---
это в точности $\theta_\gamma$ (перенесённые на $F[\alpha]$), $\gamma\in\Om$. Гомоморфизм $\theta_\gamma$ продолжается на поле
частных $F(\alpha)$ тогда и только тогда, когда $(\varphi_0h)(\gamma)\ne0$ для всех $h\ne0$. Когда $h$ пробегает ненулевые
многочлены из $F[X]$, многочлен $\varphi_0h$ пробегает все ненулевые многочлены из $\varphi_0(F)[X]$ (так как
$\varphi_0\colon F\to\varphi_0(F)$~--- изоморфизм). Значит, условие означает, что $\gamma$ трансцендентен над $\varphi_0(F)$.

(b) Если $\varphi$~--- продолжение $\varphi_0$ на $F[\alpha]$, то, применяя $\varphi$ к $\sum a_i\alpha^i=0$, получаем
$\sum\varphi_0(a_i)\varphi(\alpha)^i=0$, т.\,е.\ $\varphi(\alpha)$~--- корень $\varphi_0f$. Обратно, если $(\varphi_0f)(\gamma)=0$,
то $f\in\ker\theta_\gamma$, гомоморфизм $\theta_\gamma$ пропускается через $F[X]/(f)\cong F[\alpha]$ и даёт продолжение
$\varphi_0$ с $\alpha\mapsto\gamma$. Число различных корней $\varphi_0f$ не больше $\deg\varphi_0f=\deg f$.
\end{proof}

\par\noindent\begin{minipage}{\linewidth}
\begin{center}
\begin{tikzcd}[column sep=huge,row sep=large]
F[\alpha]\arrow[r,cnew,dashed,"\alpha\,\mapsto\,\gamma"] & \Om\\
F\arrow[u,hook]\arrow[ur,cold,"\varphi_0"'] &
\end{tikzcd}
\end{center}
\figcap{Продолжение $\varphi_0$ на $F[\alpha]$ существует ровно для тех $\gamma\in\Om$, которые являются корнями
многочлена $\varphi_0f$, а не самого $f$.}
\end{minipage}\par\medskip

\begin{example}[2.4a]
Пусть $\alpha=\sqrt[4]2\in\R$. Многочлен $X^4-2$ неприводим над $\Q$ (Эйзенштейн, $p=2$), поэтому $[\Q[\alpha]:\Q]=4$,
$[\Q[\sqrt2]:\Q]=2$ и $[\Q[\alpha]:\Q[\sqrt2]]=2$; минимальный многочлен $\alpha$ над $\Q[\sqrt2]$~--- это $f=X^2-\sqrt2$.
Возьмём $\varphi_0\colon\Q[\sqrt2]\to\CC$, $\sqrt2\mapsto-\sqrt2$ (\rf{2.1}(b)). Тогда $\varphi_0f=X^2+\sqrt2$, его корни
$\pm i\sqrt[4]2$ не вещественны. Значит, $\varphi_0$ имеет ровно два продолжения на $\Q[\alpha]$ со значениями в $\CC$
и ни одного со значениями в $\R$, хотя тождественное вложение $\Q[\sqrt2]\subset\R$ продолжается в $\R$ двумя способами
($\alpha\mapsto\pm\sqrt[4]2$). Ответ зависит не только от полей, но и от самого $\varphi_0$.
\end{example}

\subsection*{Присоединение второго корня}

Поле корня неприводимого многочлена определено однозначно с точностью до изоморфизма (\rf{1.27b}). Предложение
\rf{2.4} позволяет понять, что происходит при присоединении \emph{второго} корня.

\begin{proposition}[2.4b]
Пусть $f\in F[X]$~--- унитарный неприводимый многочлен, $(F[\alpha],\alpha)$~--- его поле корня, и
\[
f=(X-\alpha)\,g_1\cdots g_s\qquad\text{в }F[\alpha][X],\quad g_j\text{ унитарны и неприводимы}.
\]
Пусть $\Om$~--- расширение $F$, а $\alpha',\beta'\in\Om$~--- корни $f$, и $\varphi\colon F[\alpha]\to\Om$~--- $F$-гомоморфизм с
$\varphi(\alpha)=\alpha'$ (\rf{2.1}). Если $\beta'\ne\alpha'$, то $\beta'$~--- корень $\varphi g_j$ ровно для одного
многочлена $g_j$ (с точностью до повторов в списке), и существует изоморфизм
\[
F[\alpha][X]/(g_j)\xrightarrow{\ \sim\ }F[\alpha',\beta'],\qquad \alpha\mapsto\alpha',\quad x\mapsto\beta'
\]
($x$~--- класс $X$). Итак, поле $F[\alpha',\beta']$ (вместе с парой образующих) с точностью до изоморфизма определяется
тем, корнем \emph{какого} из многочленов $\varphi g_j$ является второй корень. Если все $g_j$ совпадают (например,
$s=1$), все такие поля изоморфны; если среди $g_j$ есть различные, они могут быть неизоморфны и даже иметь разные
степени над $F$.
\end{proposition}

\begin{proof}
Гомоморфизм $\varphi$~--- изоморфизм $F[\alpha]$ на подполе $F[\alpha']\subset\Om$, тождественный на $F$; применяя его к
коэффициентам, получаем в $\Om[X]$
\[
f=\varphi f=(X-\alpha')\,\varphi g_1\cdots\varphi g_s .
\]
Поэтому корень $\beta'\ne\alpha'$ является корнем некоторого $\varphi g_j$. Многочлен $\varphi g_j$ унитарен и неприводим над
$F[\alpha']$ (изоморфизм колец $F[\alpha][X]\to F[\alpha'][X]$, индуцированный $\varphi$, сохраняет неприводимость), поэтому это
минимальный многочлен $\beta'$ над $F[\alpha']$ (\rf{1.27c}); в частности, $j$ определено однозначно с точностью до
совпадения многочленов в списке. По \rf{2.4}(b), применённому к простому расширению $F[\alpha][X]/(g_j)$ поля $F[\alpha]$,
гомоморфизму $\varphi$ и корню $\beta'$ многочлена $\varphi g_j$, или непосредственно по \rf{1.27c}, отображение $F[\alpha][X]\to F[\alpha',\beta']$, $\sum c_iX^i\mapsto\sum\varphi(c_i)\beta'^i$,
сюръективно и имеет ядро $(g_j)$.
\end{proof}

\begin{example}[2.4c]
(1) $f=X^4-2$ над $\Q$. Пусть $\alpha$~--- корень, $\Q[\alpha]\cong\Q[\sqrt[4]2]\subset\R$ (\rf{2.1}(b)). В $\Q[\alpha][X]$
\[
X^4-2=(X-\alpha)(X+\alpha)(X^2+\alpha^2),
\]
и $X^2+\alpha^2$ неприводим над $\Q[\alpha]$: при вложении $\alpha\mapsto\sqrt[4]2$ он переходит в $X^2+\sqrt2$, корни
которого $\pm i\sqrt[4]2$ не лежат в $\R\supset\Q[\sqrt[4]2]$. Итак, $s=2$. Если второй корень~--- $\beta=-\alpha$, то
$\Q[\alpha,\beta]=\Q[\alpha]$ имеет степень $4$; если же $\beta$~--- корень $X^2+\alpha^2$ (в $\CC$: $\alpha=\sqrt[4]2$,
$\beta=\pm i\sqrt[4]2$), то $\Q[\alpha,\beta]=\Q[\sqrt[4]2,i]$~--- поле разложения $X^4-2$ (т.\,е.\ поле, порождённое
\emph{всеми} корнями $X^4-2$; см.\ раздел <<Поля разложения>>) степени $8$. Поля разных степеней
не изоморфны, так что слова <<поле, порождённое над $\Q$ двумя корнями $X^4-2$>> смысла не имеют.

(2) $f=X^3-2$ над $\Q$. Здесь $X^3-2=(X-\alpha)(X^2+\alpha X+\alpha^2)$ в $\Q[\alpha][X]$, и квадратичный множитель
неприводим (его корни $\omega\sqrt[3]2,\omega^2\sqrt[3]2$ при $\alpha=\sqrt[3]2$ не вещественны). Здесь $s=1$: любые два
различных корня порождают поле разложения $X^3-2$ степени $6$.
\end{example}

\par\noindent\begin{minipage}{\linewidth}
\begin{center}
\begin{tikzpicture}[x=1cm,y=1cm]
\node (Q) at (0,0) {$\Q$};
\node (A) at (0,1.5) {$\Q[\alpha]=\Q[\alpha,-\alpha]$};
\node (B) at (0,3.2) {$\Q[\alpha,i\alpha]$};
\draw (Q) -- node[left,font=\footnotesize]{$4$} (A);
\draw[cnew,thick] (A) -- node[left,font=\footnotesize]{$2$} (B);
\node[font=\footnotesize,align=left,cold] at (4.6,1.5) {$\beta=-\alpha$: корень множителя $X+\alpha$};
\node[font=\footnotesize,align=left,cnew] at (4.6,3.2) {$\beta=\pm i\alpha$: корень множителя $X^2+\alpha^2$};
\end{tikzpicture}
\end{center}
\figcap{Присоединение второго корня $X^4-2$: результат зависит от того, корнем какого неприводимого множителя над
$\Q[\alpha]$ является второй корень.}
\end{minipage}\par\medskip

\begin{aside}[2.4d]
На языке теории Галуа (глава~3) картина такова. Пусть $f$ неприводим и сепарабелен (\rf{2.22}), $L$~--- его поле разложения,
$G=\Gal(L/F)$ и $H=\Gal(L/F[\alpha])$~--- стабилизатор корня $\alpha$. Тогда неприводимые множители $f$ над $F[\alpha]$
соответствуют орбитам $H$ на корнях $f$: множителю отвечает произведение $\prod_{\gamma\in\mathcal O}(X-\gamma)$ по
орбите $\mathcal O$ (его коэффициенты неподвижны относительно $H$, т.\,е.\ лежат в $F[\alpha]$, а любой его делитель над
$F[\alpha]$ вместе с корнем $\gamma$ имеет корнями все $h\gamma$, $h\in H$). Поэтому $[F[\alpha,\beta]:F[\alpha]]=|H\beta|$.
Если $H$ транзитивна на остальных корнях, т.\,е.\ $G$ действует на корнях \emph{дважды транзитивно} (например,
$G=S_n$, а также $G=A_n$ при $n\ge4$), то множитель $g_1$ один, и <<присоединение двух корней>> определено однозначно (с точностью до изоморфизма).
Обратное неверно: разные орбиты могут давать изоморфные поля (см.\ пример $X^5-5X+12$ ниже). Для $X^4-2$ группа
$G$~--- группа симметрий квадрата с вершинами $\pm\sqrt[4]2,\pm i\sqrt[4]2$, стабилизатор $\sqrt[4]2$ состоит из
тождества и комплексного сопряжения, и его орбиты $\{\alpha\}$, $\{-\alpha\}$, $\{\pm i\alpha\}$ дают множители из
\rf{2.4c}(1).

Неоднозначность бывает и тогда, когда множители имеют одинаковую степень. Пусть $\rho_0=\rho,\rho_1,\rho_2$~--- корни
$X^3-X-1$ в $\CC$, где $\rho$ вещественный, $\theta_{ij}=\omega^i\sqrt[3]2+\rho_j$ ($i,j=0,1,2$), $\theta=\theta_{00}=\sqrt[3]2+\rho$ и
\[
f=\prod_{i,j=0}^{2}(X-\theta_{ij})=\prod_{j=0}^{2}\bigl((X-\rho_j)^3-2\bigr)=X^9-3X^7-9X^6+3X^5-28X^3-9X^2+27X-25
\]
(коэффициенты~--- симметрические многочлены от $\rho_0,\rho_1,\rho_2$ с рациональными коэффициентами и потому
рациональны). Поля разложения $L_1=\Q[\sqrt[3]2,\omega]$ и $L_2=\Q[\rho_0,\rho_1,\rho_2]$ многочленов $X^3-2$ и $X^3-X-1$
(оба неприводимы: у них нет рациональных корней) имеют степень $6$ и группу Галуа $S_3$, так как их дискриминанты
$-108=-3\cdot6^2$ и $-23$~--- не квадраты (\rf{2.6}); их единственные квадратичные подполя (в $S_3$ одна подгруппа
индекса $2$)~--- это $\Q[\sqrt{-3}]$ и $\Q[\sqrt{-23}]$ (\rf{4.2}). Эти квадратичные поля различны, так как
$(-3)(-23)=69$~--- не квадрат в $\Q$. Пересечение $L_1\cap L_2$ нормально над $\Q$ (\rf{3.21}), а нормальные подполя
$S_3$-расширения~--- это $\Q$, квадратичное подполе и всё поле (нормальные подгруппы $S_3$~--- это $1$, $A_3$, $S_3$);
так как квадратичные подполя $L_1$ и $L_2$ различны (и потому $L_1\ne L_2$), имеем $L_1\cap L_2=\Q$. Поэтому по
\rf{3.21a} композит $L=L_1L_2$ имеет степень $36$ и группу Галуа $G=S_3\times S_3$: первый сомножитель переставляет
$\omega^i\sqrt[3]2$, второй~--- $\rho_j$, т.\,е.\ $G$ действует на таблице корней $\theta_{ij}$ перестановками строк и
столбцов.

\emph{Корни $\theta_{ij}$ различны.} Если $\theta_{ij}=\theta_{kl}$, то $d=\rho_l-\rho_j=\omega^i\sqrt[3]2-\omega^k\sqrt[3]2\in
L_1\cap L_2=\Q$. Пусть $j\ne l$. Циклическая перестановка $\rho_0\to\rho_1\to\rho_2\to\rho_0$ лежит в $\Gal(L_2/\Q)=S_3$ и
оставляет $d$ на месте, поэтому $\rho_{l+1}-\rho_{j+1}=\rho_{l+2}-\rho_{j+2}=d$ (индексы по модулю $3$); сложение трёх
равенств даёт $3d=0$, т.\,е.\ $\rho_l=\rho_j$~--- противоречие. Значит, $j=l$, $d=0$ и $i=k$. Элемент $G$, оставляющий на
месте все клетки таблицы, тождествен, поэтому поле, порождённое всеми $\theta_{ij}$, равно $L$ (основная теорема теории
Галуа), т.\,е.\ $L$~--- поле разложения $f$.

\emph{Неприводимость и разложение над $\Q[\theta]$.} Группа $G$ транзитивна на клетках, поэтому минимальный многочлен
$\theta$ имеет корнями все девять $\theta_{ij}$, и $f$ неприводим. Стабилизатор клетки $(0,0)$~--- это
$H=\Gal(L/\Q[\theta])=S_2\times S_2$; он оставляет на месте $\sqrt[3]2$ и $\rho$, так что $\Q[\theta]=\Q[\sqrt[3]2,\rho]$
(степень $36/4=9$). Орбиты $H$ на клетках~--- это $\{(0,0)\}$, $\{(1,0),(2,0)\}$, $\{(0,1),(0,2)\}$ и четыре
оставшиеся клетки; по сказанному выше над $\Q[\theta]$ многочлен $f$ раскладывается на неприводимые множители степеней
$1,2,2,4$. Квадратичные множители можно выписать явно: так как $\omega\sqrt[3]2+\omega^2\sqrt[3]2=-\sqrt[3]2$,
$\omega\sqrt[3]2\cdot\omega^2\sqrt[3]2=\sqrt[3]4$, $\rho_1+\rho_2=-\rho$ и $\rho_1\rho_2=1/\rho=\rho^2-1$ (из $\rho^3=\rho+1$),
\begin{gather*}
(X-\theta_{10})(X-\theta_{20})=(X-\rho)^2+\sqrt[3]2\,(X-\rho)+\sqrt[3]4,\\
(X-\theta_{01})(X-\theta_{02})=(X-\sqrt[3]2)^2+\rho\,(X-\sqrt[3]2)+\rho^2-1 .
\end{gather*}
Их неприводимость над $\Q[\theta]=\Q[\sqrt[3]2,\rho]$ видна и напрямую. Если бы первый множитель был приводим, то
$\theta_{10}\in\Q[\theta]$, а с ним и $\omega=(\theta_{10}-\rho)/\sqrt[3]2$ степени $2$ над $\Q$; если бы второй, то
$\rho_1=\theta_{01}-\sqrt[3]2\in\Q[\theta]$, а значит, и $L_2=\Q[\rho,\rho_1]$ степени $6$. Но ни $2$, ни $6$ не делят
$9=[\Q[\theta]:\Q]$ (\rf{1.20}).

\emph{Три поля.} Присоединение корня $\theta_{10}=\omega\sqrt[3]2+\rho$ первого квадратичного множителя даёт поле
неподвижных точек стабилизатора клеток $(0,0)$ и $(1,0)$, т.\,е.\ подгруппы $1\times S_2$; это $L_1[\rho]=\Q[\sqrt[3]2,\rho,\omega]$.
Корень $\theta_{01}=\sqrt[3]2+\rho_1$ второго даёт поле неподвижных точек подгруппы $S_2\times1$, т.\,е.\
$L_2[\sqrt[3]2]=\Q[\sqrt[3]2,\rho,\rho_1]$. Корень $\theta_{11}$ множителя степени $4$ даёт всё $L$ (стабилизатор клеток
$(0,0)$ и $(1,1)$ тривиален). Первые два поля имеют степень $36/2=18$, но не изоморфны: в первом $-3$ является
квадратом ($\sqrt{-3}=2\omega+1$), а во втором нет~--- иначе оно содержало бы $\omega$, т.\,е.\ и $L_1$, и $L_2$, а
значит, и $L$ степени $36>18$.

Даже при самой большой группе Галуа <<присоединение двух корней>> может быть неоднозначным. Пусть $x_1,\dots,x_5$~---
корни $X^5-X-1$ в $\CC$; группа Галуа этого многочлена~--- $S_5$ (\rf{4.30}). Многочлен $P=\prod_{i<j}(X-x_i-x_j)$
степени $10$ имеет рациональные коэффициенты (это симметрические многочлены от $x_i$, как и выше):
$P=X^{10}+3X^6+11X^5-4X^2+4X-1$. Все числа $x_i+x_j$ различны (с точностью до нумерации корней достаточно двух проверок): из $x_1+x_2=x_1+x_3$
следует $x_2=x_3$, а из $x_1+x_2=x_3+x_4$ после применения транспозиции $(2\,5)\in S_5$ получаем
$x_1+x_5=x_3+x_4=x_1+x_2$, т.\,е.\ $x_5=x_2$. Поэтому стабилизатор
корня $x_1+x_2$~--- стабилизатор пары $\{1,2\}$, и на остальных парах у него две орбиты: три пары, не пересекающиеся с
$\{1,2\}$, и шесть пар с одним общим элементом. Значит, над $\Q[x_1+x_2]$ многочлен $P$ раскладывается на неприводимые
множители степеней $1,3,6$, и второй корень даёт поля степеней $120/4=30$ и $120/2=60$ над $\Q$.

Наоборот, для $f=X^5-5X+12$ над $\Q$ группа Галуа $G$~--- группа диэдра $D_5$ порядка $10$, действующая на пяти корнях
как на вершинах правильного пятиугольника (ср.\ упражнение 4-7). Проверим это. Многочлен $f$ неприводим уже по модулю
$7$ (\rf{1.18}), поэтому для его корня $\alpha$ степень $5=[\Q[\alpha]:\Q]$ делит $|G|$. Дискриминант $f$ равен
$2^{12}5^6=8000^2$, так что $G\subset A_5$ (\rf{4.2}). Вещественный корень у $f$ один ($f'=5(X^4-1)$, а значения в
точках локального максимума и минимума $f(-1)=16$, $f(1)=8$ положительны), поэтому комплексное сопряжение даёт в $G$
элемент порядка $2$. Наконец, перемножение с учётом $\alpha^5=5\alpha-12$ показывает, что в $\Q[\alpha][X]$
\[
f=(X-\alpha)\Bigl(X^2-\frac{u-3\alpha-4}4\,X-\frac{u+5\alpha-8}4\Bigr)\Bigl(X^2+\frac{u+\alpha-4}4\,X-\frac{\alpha^3+\alpha+2}2\Bigr),
\]
где $u=\alpha^4+\alpha^3+\alpha^2$, так что поле разложения получается из $\Q[\alpha]$ присоединением корней двух квадратных трёхчленов, и $|G|\le5\cdot2\cdot2=20$.
Итак, $|G|$ делится на $10$, делит $60$ и не больше $20$; подгрупп порядка $20$ (индекса $3$) в простой группе $A_5$ нет,
поэтому $|G|=10$. Группа порядка $10$ в $S_5$ не циклическая (элементов порядка $10$ в $S_5$ нет), т.\,е.\ это $D_5$.
Стабилизатор корня $\alpha$~--- отражение; его орбиты на корнях~--- $\{\alpha\}$ и две пары, т.\,е.\ выписанные
квадратные трёхчлены неприводимы над $\Q[\alpha]$, и разложение имеет степени $1,2,2$. Но элемент $D_5$,
оставляющий на месте две вершины, тождествен, поэтому любые два корня порождают всё поле разложения (степени $10$):
оба квадратичных множителя дают одно и то же поле.
\end{aside}
